\section{INTRODUCTION}

This manual describes the \textbf{automation module} of \textbf{FAST} (Forecast Automation System for
Telescopes), the component that runs a single, once-per-night Meso-NH / Astro-Meso-NH atmospheric and
optical-turbulence forecast for a ground-based telescope site. It automates the full chain: waiting for
and validating the initialization files, running the mesoscale simulation, post-processing the raw
model output into astroclimatic parameters, generating diagnostic plots, and delivering the results to
a remote server.

This document describes the automation \emph{generically}, i.e.\ independently of the specific
telescope, site or forecasting project it is deployed for. FAST has been deployed, in slightly
different configurations, for more than one ground-based observatory; this manual covers the common
core that any such deployment shares, not the site-specific choices (domain geometry, delivered file
formats, calibration constants) that a given installation layers on top of it via its own
\texttt{conf.d/*.var} files. Where a parameter table below shows a \texttt{[current: ...]} value, that
value is simply whatever this particular installation's configuration files currently contain, given as
a worked example -- not a universal default.

This document covers only the \emph{automation} part of FAST, i.e.\ the component that produces the
once-per-night Meso-NH forecast. A separate, complementary \emph{short-timescale} component (an
autoregressive or machine-learning correction of the forecast using real-time telemetry, updated every
few minutes) may run alongside it on some deployments; it is a distinct program with its own
configuration and is documented separately. Section~\ref{sec:usage} describes the point at which the
automation hands data off to it, where applicable.

The user of this manual is assumed to be familiar with the basic operation of Meso-NH, with the Linux
shell, and with basic bash scripting. A working installation of Meso-NH and of a local OpenMPI build is
assumed to already be in place; Section~\ref{sec:deps-install} lists the remaining dependencies and the
steps needed to install and configure the automation code on top of them.

\textbf{Scope of this revision.} This is a \emph{single-configuration} edition of the manual: it
documents one self-contained instance of the automation (one \texttt{bin/} tree, one set of
\texttt{conf.d/*.var} files, one nightly \texttt{run.sh}), started directly (by cron or manually) with
no external orchestration layer coordinating multiple instances or deciding which one's output to keep.
Deployments that do chain several configurations together (for example, running a second, different-
resolution configuration after the first, or driving several instances from a shared top-level script)
require that coordination layer to be documented separately, on top of this manual; nothing here
describes it.

\textbf{A note on scope.} This manual documents what the automation does and how it is configured; it
assumes solid Linux system administration and shell-scripting skills, and it is not a substitute for
them. In an operational deployment, reading the source code directly -- not just this manual -- is
eventually unavoidable: some situation not covered here will arise, and resolving it will require
understanding the scripts themselves. No manual can substitute for the operational experience needed to
run an unattended automated system in production; that experience is built on the job, by encountering
and working through real problems as they occur.
